\section{Experiments}

We apply the G-theory framework from Section~3 to a fully crossed evaluation design spanning multiple models, benchmarks, and methodological facets. All experiments use open-weight 7--9B parameter models evaluated with vLLM on a single A100-80GB GPU. Binary correctness (1/0) serves as the primary outcome measure throughout.

\subsection{Single-Model Baseline}

We first establish the variance structure within a single model (Llama-3.1-8B-Instruct) on MMLU to isolate methodological facet contributions before introducing cross-model variation.

\paragraph{Design.} The G-study crosses six facets: 200 items $\times$ 3 numerical precisions (FP16, BF16, FP32) $\times$ 6 prompt templates $\times$ 6 random seeds $\times$ 3 temperature levels (0.0, 0.3, 0.7) $\times$ 4 answer orderings, yielding 239,200 binary observations. Variance components are estimated via Henderson Method~I with bootstrap confidence intervals (1,000 resamples).

\paragraph{Results.} Under binary scoring, item difficulty (\texttt{item\_id}) accounts for 61.1\% of total variance, with the residual term contributing 19.5\%. All main facet effects are negligible ($<$1\% each): prompt template 0.3\%, temperature 0.2\%, precision 0.1\%, seed $<$0.1\%, ordering $<$0.1\%. The interesting structure appears in item-conditional interactions: precision$\times$item contributes 7.2\%, temperature$\times$item 6.9\%, and prompt$\times$item 5.7\%. This pattern indicates that methodological choices do not shift performance uniformly across items; instead, specific items are differentially sensitive to specific facets.

The resulting generalizability coefficient is $G_\text{item} = 0.918$, indicating that 200 items under multi-facet conditions produce highly reliable rankings. A D-study projection shows $G \geq 0.80$ is achievable with as few as 72 items in this multi-facet design.

For text-exact-match scoring (requiring verbatim string matches), the variance structure shifts dramatically: temperature dominates at 79.7\% of variance, and $G_\text{item}$ collapses to 0.007. This metric-conditional divergence motivates our exclusive use of binary correctness for the remaining analyses.

\subsection{Cross-Model Variance}

We extend the single-model design to four models: Llama-3.1-8B-Instruct, Gemma-2-9B-it, Mistral-7B-Instruct-v0.3, and Qwen3-8B. A balanced BF16 subset (345,600 records) serves as the primary analysis sample, chosen because BF16 is the standard inference precision for these architectures.

\paragraph{Results.} Table~\ref{tab:cross_model} presents the variance decomposition. Introducing model as a facet reshapes the structure: item difficulty drops from 61.1\% to 42.1\% [36.0--48.8\%], while model$\times$item interaction emerges as the second-largest component at 33.6\% [28.4--37.9\%]. The model main effect contributes only 0.7\%, confirming that models differ not in overall ability level but in \emph{which items} they find difficult.

\begin{table}[t]
\centering
\caption{Variance decomposition for the cross-model MMLU design (4 models, BF16, 345,600 records). Components estimated via Henderson Method~I with 1,000-resample bootstrap CIs. $G_\text{item} = 0.820$ [0.782--0.861].}
\label{tab:cross_model}
\small
\begin{tabular}{lrl}
\toprule
Component & Var.\% & 95\% CI \\
\midrule
\textbf{item\_id} & \textbf{42.1} & [36.0--48.8] \\
\underline{model$\times$item} & \underline{33.6} & [28.4--37.9] \\
residual & 19.5 & [17.2--21.8] \\
precision$\times$item & 1.3 & [0.8--1.9] \\
model & 0.7 & [0.2--1.5] \\
prompt$\times$item & 0.6 & [0.3--1.0] \\
temperature$\times$item & 0.6 & [0.3--1.0] \\
seed$\times$item & 0.4 & [0.1--0.8] \\
ordering$\times$item & 0.3 & [0.1--0.6] \\
all others & $<$0.1 & \\
\bottomrule
\end{tabular}
\end{table}

The per-model decomposition of interaction magnitude reveals meaningful asymmetry: Mistral concentrates 29.3\% of the total model$\times$item variance, followed by Qwen (27.1\%), Gemma (24.0\%), and Llama (19.7\%). This ordering suggests that instruction-tuning strategies and pretraining data composition produce model-specific item difficulty profiles.

\paragraph{D-study validation.} To test whether G-theory projections generalize beyond the training sample, we estimate a D-study on 3-model subsets and predict the 4-model coefficient. The predicted $G = 0.816$ for $n_\text{model} = 4$ closely matches the observed $G = 0.820$ ($\Delta = 0.004$). Forward projections indicate $G = 0.848$ at 5 models, $G = 0.893$ at 8 models, and $G \geq 0.90$ requires approximately 10 models.

\subsection{Format-Conditional Variance Structure}
\label{sec:format_conditional}

The preceding analyses use only MMLU, a multiple-choice benchmark. We now test whether the item-dominated variance structure generalizes across evaluation formats by extending to four benchmarks spanning two response formats.

\paragraph{Design.} The expanded design (exp-002) crosses 4 models $\times$ 4 benchmarks $\times$ BF16 $\times$ 2 temperatures (0.0, 0.7) $\times$ 6 prompts $\times$ 6 seeds $\times$ 4 orderings $\times$ 200 items per benchmark, totaling 921,600 records. The benchmarks include three multiple-choice tasks (MMLU, ARC-Challenge, HellaSwag) and one free-form generation task (GSM8K). For free-form tasks, the ordering facet permutes few-shot exemplar sequences rather than answer options.

% [PLACEHOLDER: expand to 8 models × 5 benchmarks (+ MATH) when data ready]

\paragraph{Results.} Table~\ref{tab:cross_benchmark} reveals a qualitative structural shift between formats. For all three multiple-choice benchmarks, item difficulty is the dominant component (MMLU 42\%, ARC 49\%, HellaSwag 45\%), and model$\times$item interaction is the second largest (32--38\%). This structure is consistent with the cross-model analysis in Section~4.2.

GSM8K exhibits a fundamentally different variance profile. Item difficulty accounts for only 12\% of variance. Instead, model (15\%), prompt template (12\%), and model$\times$prompt interaction (9\%) become the leading components. The residual term absorbs 35\% of variance, indicating substantial unexplained variation from seed and temperature effects in free-form generation.

\begin{table}[t]
\centering
\caption{Format-conditional variance structure across four benchmarks (exp-002, 921,600 records). Top-5 components plus residual per benchmark. MC = multiple-choice, FF = free-form.}
\label{tab:cross_benchmark}
\small
\begin{tabular}{l rrrr}
\toprule
 & \multicolumn{3}{c}{MC} & FF \\
\cmidrule(lr){2-4} \cmidrule(lr){5-5}
Component & MMLU & ARC & Hella. & GSM8K \\
\midrule
\textbf{item\_id} & \textbf{42} & \textbf{49} & \textbf{45} & 12 \\
model$\times$item & 34 & 32 & 38 & 8 \\
\textbf{model} & 1 & 1 & 1 & \textbf{15} \\
prompt & $<$1 & $<$1 & $<$1 & 12 \\
model$\times$prompt & $<$1 & $<$1 & $<$1 & 9 \\
residual & 19 & 15 & 13 & 35 \\
\bottomrule
\end{tabular}
% [PLACEHOLDER: add MATH column + 8-model rows when data ready]
\end{table}

This dichotomy is not a continuous gradient. The gap between the lowest MC item dominance (MMLU, 42\%) and GSM8K (12\%) represents a qualitative shift in what drives evaluation variance. In multiple-choice settings, the item pool determines reliability; in free-form settings, prompt design and model selection dominate. Figure~\ref{fig:cross_benchmark} visualizes this structural contrast.

\begin{figure}[t]
\centering
\includegraphics[width=\columnwidth]{fig1_cross_benchmark}
\caption{Variance decomposition by benchmark format. Multiple-choice benchmarks (MMLU, ARC, HellaSwag) share an item-dominated structure, while free-form GSM8K distributes variance across model, prompt, and interaction components.}
\label{fig:cross_benchmark}
\end{figure}

\subsection{Precision Separability}

Prior work has identified numerical precision as a potential source of evaluation variance \cite{yuan2025layercast}. Our multi-facet design allows us to isolate the precision contribution while controlling for all other facets simultaneously.

Under binary scoring with BF16 and FP16, the mean accuracy difference is 0.23 percentage points. After excluding FP32, the precision main effect drops to 0.0006\% of total variance. The FP32 anomaly traces to a vLLM implementation detail (\texttt{enforce\_eager} mode required for FP32 inference), not a fundamental precision effect.

Precision does not interact with other methodological facets under binary scoring: all precision$\times$methodology interactions contribute $<$0.1\% of variance. Under text-exact-match scoring, however, a precision$\times$temperature interaction emerges at 1.3\%, consistent with the metric-conditional pattern observed in Section~4.1. The practical implication is clear: for standard binary evaluation at BF16/FP16, numerical precision is a negligible variance source that need not be controlled. This finding is specific to single-GPU inference and does not address multi-GPU tensor parallelism, where numerical non-associativity introduces additional variance sources \cite{shanmugavelu2024floating}.

\subsection{Design-Dependent Reproducibility Budgets}

The previous sections establish that item selection and model$\times$item interaction dominate evaluation variance. A natural question follows: how many items does a given evaluation design need to achieve a target reliability? We address this through D-study projections that translate variance components into concrete item budgets.

Table~\ref{tab:dstudy} compares four evaluation designs at the $G \geq 0.80$ threshold. The multi-facet design (432 conditions per item: 3 precisions $\times$ 6 prompts $\times$ 6 seeds $\times$ 4 orderings) achieves $G = 0.918$ at 200 items and requires only 72 items for $G \geq 0.80$. A single-condition design (one fixed prompt, seed, ordering, and precision per item) reaches only $G = 0.559$ at 200 items and requires 631 items for the same threshold. The ratio between these budgets is 8.8$\times$, and this ratio remains stable across subset sizes from 50 to 200 items (CV = 1.9\%).

\begin{table}[t]
\centering
\caption{D-study reproducibility budgets under different evaluation designs. The multi-facet design requires $8.8\times$ fewer items than single-condition evaluation to reach $G \geq 0.80$.}
\label{tab:dstudy}
\small
\begin{tabular}{lcrrr}
\toprule
Design & Cond./item & $n_{G \geq 0.80}$ & $G_{200}$ & $G_{116}$ \\
\midrule
Multi-facet & 432 & \textbf{72} & \textbf{0.918} & \textbf{0.882} \\
Partial (4 cond.) & 4 & 189 & 0.810 & 0.721 \\
Single-condition & 1 & 631 & 0.559 & 0.424 \\
\midrule
IRT baseline & \multicolumn{1}{c}{---} & 116 & --- & 0.424$^\dagger$ \\
\bottomrule
\end{tabular}

\vspace{2pt}
{\scriptsize $^\dagger$\,PSN-IRT item count of 116 \cite{zhou2026irt} evaluated under single-condition protocol yields $G = 0.424$. See text for interpretation.}
\end{table}

The comparison with IRT-based efficient benchmarking is instructive. \citet{zhou2026irt} use IRT to select 116 informative items from MMLU, optimizing for item discrimination. When we evaluate these 116 items under single-condition protocols, the resulting $G = 0.424$ falls well below the reliability threshold. This is not a failure of IRT item selection; rather, it reflects that IRT optimizes a different objective (latent trait estimation) that does not account for methodological facet variance. Under the multi-facet design, 116 items yield $G = 0.882$, confirming that the reliability gap is driven by evaluation design, not item quality.

Figure~\ref{fig:dstudy} plots the D-study projection curves for each design. The multi-facet curve saturates rapidly (diminishing returns beyond $\sim$150 items), while the single-condition curve rises slowly and does not reach $G = 0.80$ within typical benchmark sizes.

\begin{figure}[t]
\centering
\includegraphics[width=\columnwidth]{fig3_dstudy_curves}
\caption{D-study projection curves for three evaluation designs. The multi-facet design ($8.8\times$ fewer items) reaches $G \geq 0.80$ at 72 items; the single-condition design requires 631 items. The shaded band shows 95\% bootstrap CIs.}
\label{fig:dstudy}
\end{figure}

\subsection{Robustness Checks}

We conduct six robustness analyses to verify that the reported variance structures are not artifacts of specific analytical choices. \textbf{Item resampling} (exp-001b): three independent 200-item samples drawn from the full MMLU pool yield item difficulty variance shares of 59.8\%, 61.4\%, and 62.1\% (CV = 2.7\%) and $G_\text{item}$ values with CV = 0.5\%. \textbf{Estimator robustness}: REML and Henderson Method~I differ by $<$0.5 percentage points on all components. \textbf{CI method}: bootstrap and jackknife confidence intervals agree closely (width ratio 0.96--1.04 across components). \textbf{Facet specification}: treating facets as fixed vs.\ random changes $G_\text{item}$ by only 0.37 percentage points. \textbf{Leave-one-model-out (LOMO)}: all three-model subsets preserve the item + model$\times$item dominance pattern (item\_id 39--45\%, model$\times$item 32--37\%, $G_\text{item}$ 0.747--0.793). \textbf{Null calibration}: a Bernoulli simulation with matched item difficulties produces a variance ratio of 1.0008 relative to observed data, confirming that the decomposition does not introduce systematic bias on binary outcomes.

% [PLACEHOLDER: §4.7 Extended Results — 8-model × 5-benchmark expansion]
